\section{Probabilistic Dependency Risk}
\label{sec:dependency}

The companion paper \citep{lasser2026horizon} bounds the worst-case
contraction of single-substrate versus cross-substrate populations
through a minimax analysis (Proposition~4). That bound is a
one-time penalty on the value function: useful for showing
diversification is structurally preferred, but not directly usable
inside a discounted trajectory because it does not attach a per-step
cost. This section converts the minimax bound into an expected
per-step dependency-risk cost, then folds that cost into the
discounted value function so diversification enters the horizon
calculation as a first-class term rather than a static penalty.

The section has three parts. First, a probabilistic substrate threat
model (Section~\ref{sec:threat_model_below}) that decomposes
dependency risk into event probability, target selection, and cascade
effects. Second, an expected per-step risk comparison
(Proposition~\ref{prop:expected_risk}) that sharpens the minimax
bound into an \emph{adversarial balance condition} $c > \max_i (f_i
\cdot c_i)$ --- a single inequality that gates when diversification
strictly beats concentration and that explicitly admits contagion
dynamics (pandemics, cascading failures, memetic attacks) as the
motivating threat class. Third, a risk-adjusted trajectory analysis
(Section~\ref{sec:risk_adjusted}) showing that the per-step cost
compounds at the same discount rate as growth, so the
diversification-versus-concentration decision becomes a direct
comparison of two sums in the value function.

\subsection{Substrate Threat Model}
\label{sec:threat_model_below}

\begin{definition}[Substrate Threat Model]
\label{def:threat_model}
A \emph{substrate threat model} is a distribution over adversarial events,
    parametrized by:
\begin{itemize}
\item $p_{\mathrm{adv}}$: per-step probability that a substrate-targeting
    event occurs.
\item $q_i = p(\mathcal{S}_i \text{ targeted} \mid \text{event occurs})$:
    the conditional probability that substrate class $\mathcal{S}_i$ is
    targeted, given that an event occurs. For an uninformed adversary,
    $q_i = 1/m$ (uniform); for an optimal adversary,
    $q_i = \mathbf{1}[i = \arg\max_j (|\Delta \volL(\mathcal{S}_j)|
    \cdot c_j)]$, i.e., the adversary targets the substrate with
    maximal \emph{effective damage} (volume contribution times
    cascade fraction).
\item $c_i \in [0, 1]$: the \emph{cascade fraction} for substrate
    $\mathcal{S}_i$, defined as the expected fraction of
    $\mathcal{S}_i$'s total volume lost when a substrate-targeting
    event hits $\mathcal{S}_i$. ``Cascade'' refers to the propagation
    dynamics --- within-substrate impairment spreading from the
    initial strike through the SCM's structural equations --- and
    ``fraction'' constrains $c_i$ to $[0, 1]$ as a proportion of
    $\mathcal{S}_i$'s volume. The cascade fraction aggregates the
    initial impairment and all subsequent within-substrate
    propagation: if a direct attack impairs 30\% of the substrate
    and triggers a contagion that takes out another 40\%, then
    $c_i = 0.7$. The homogeneous single-substrate baseline uses a
    single $c$ in place of the per-substrate $c_i$.
\end{itemize}
In this threat model $|\Delta \volL(\mathcal{S}_i)|$ denotes
    substrate $\mathcal{S}_i$'s \emph{total contribution to $\volL(G)$}
    --- i.e., the maximum possible contraction if $\mathcal{S}_i$
    were fully eliminated --- not the expected contraction from a
    targeting event. The expected contraction is
    $|\Delta \volL(\mathcal{S}_i)| \cdot c_i$: substrate size times
    cascade fraction, with the two quantities varied independently.
    This convention matches Proposition~4 of \citet{lasser2026horizon},
    where $|\Delta \volL(\mathcal{S}_i)|$ is the minimax-bound
    substrate size, and prevents double-counting when the cascade
    fraction is applied as a multiplier.
\end{definition}

The threat model decomposes dependency risk into three components:
occurrence probability, target selection, and cascading effects. The
minimax analysis of \citet{lasser2026horizon} assumes $p_{\mathrm{adv}} > 0$
(non-negligible threat), $q_i$ chosen adversarially (worst case), and full
cascade within the targeted substrate (all agents on that substrate are
impaired). The probabilistic model relaxes each of these: $p_{\mathrm{adv}}$
is estimated from the actor's world model, $q_i$ reflects the actor's
beliefs about threat targeting, and cascading effects are modeled through
the SCM's structural equations.

\subsection{Expected Per-Step Dependency Risk}

\paragraph{Notation update from the companion paper.} The companion
paper \citep{lasser2026horizon} defines $\Risk_{\mathrm{dep}}$ as a
discounted pathway sum over the full planning horizon. In this
section, $\Risk_{\mathrm{dep}}^{\mathrm{single}}$ and
$\Risk_{\mathrm{dep}}^{\mathrm{mixed}}$ denote the \emph{per-step
expected dependency-risk cost} under the substrate threat model
(Definition~\ref{def:threat_model}): a single-step expected
contraction, not a discounted sum. The per-step form is the input to
the risk-adjusted trajectory comparison
(Section~\ref{sec:risk_adjusted}), which recovers a discounted sum
by standard geometric-series aggregation.

\begin{proposition}[Expected Risk Comparison]
\label{prop:expected_risk}
Under the substrate threat model
    (Definition~\ref{def:threat_model}), the expected per-step dependency
    risk for each strategy is:
\begin{align}
\label{eq:expected_risk_single}
\mathbb{E}[\Risk_{\mathrm{dep}}^{\mathrm{single}}]
    &= p_{\mathrm{adv}} \cdot \volL(G) \cdot c \\
\label{eq:expected_risk_mixed}
\mathbb{E}[\Risk_{\mathrm{dep}}^{\mathrm{mixed}}]
    &= p_{\mathrm{adv}} \cdot \sum_i q_i \cdot
    |\Delta \volL(\mathcal{S}_i)| \cdot c_i
\end{align}
where $c$ and $c_i$ are the cascade fractions from
    Definition~\ref{def:threat_model}. For a single-substrate population
    ($m = 1$), any substrate-targeting event affects a fraction $c$ of
    the entire population. For a cross-substrate population
    ($m \geq 2$):
\begin{equation}
\label{eq:risk_gap}
\mathbb{E}[\Risk_{\mathrm{dep}}^{\mathrm{single}}] -
    \mathbb{E}[\Risk_{\mathrm{dep}}^{\mathrm{mixed}}]
    = p_{\mathrm{adv}} \cdot \left(\volL(G) \cdot c
    - \sum_i q_i \cdot |\Delta \volL(\mathcal{S}_i)| \cdot c_i\right)
\end{equation}
Under \textbf{worst-case adversarial targeting}
    ($q_i = \mathbf{1}[i = \arg\max_j (|\Delta \volL(\mathcal{S}_j)|
    \cdot c_j)]$), this gap is positive whenever $p_{\mathrm{adv}} > 0$,
    the population satisfies substrate capability non-redundancy
    (Proposition~4 of \citet{lasser2026horizon}), and the
    \textbf{adversarial balance condition}
    $c > \max_i (f_i \cdot c_i)$ holds, where
    $f_i = |\Delta \volL(\mathcal{S}_i)| / \volL(G)$ is substrate
    $\mathcal{S}_i$'s share of total volume. The condition is both
    sufficient and necessary under worst-case targeting; for
    non-adversarial $q_i$ (e.g., uniform), the condition is sufficient
    but not necessary.
\end{proposition}

\begin{proof}
Under an adversarial target-selection model, $q_i = \mathbf{1}[i = i^*]$
where $i^* = \arg\max_j (|\Delta \volL(\mathcal{S}_j)| \cdot c_j)$.
Substituting into Equation~\eqref{eq:expected_risk_mixed} gives
\[
    \mathbb{E}[\Risk_{\mathrm{dep}}^{\mathrm{mixed}}]
    = p_{\mathrm{adv}} \cdot \max_i \bigl(|\Delta \volL(\mathcal{S}_i)|
        \cdot c_i\bigr)
    = p_{\mathrm{adv}} \cdot \volL(G) \cdot \max_i (f_i \cdot c_i),
\]
while $\mathbb{E}[\Risk_{\mathrm{dep}}^{\mathrm{single}}] =
p_{\mathrm{adv}} \cdot \volL(G) \cdot c$. The gap is positive iff
$c > \max_i (f_i \cdot c_i)$. Substrate capability non-redundancy
(Proposition~4 of \citet{lasser2026horizon}) ensures $f_i < 1$ strictly
for every $i$, so the condition is non-vacuous for any $c > 0$: some
non-trivial region of $(f_i, c_i)$ space satisfies it and some does
not, and the proposition is a sharp characterization of which region.
\end{proof}

\paragraph{On contagion dynamics.} The adversarial balance condition
admits contagions by construction, which is essential because
contagion-driven threats --- pandemics, cascading infrastructure
failures, financial contagion, memetic attacks --- are precisely the
class where per-substrate cascade fractions can approach saturation
($c_i \to 1$) on the targeted substrate. Smaller substrates are expected to
reach sigmoid saturation faster than larger ones because contagion
propagation time scales sub-linearly with population size in
standard epidemiological models \citep{anderson1991infectious}, so
$c_i > c$ is a plausible assumption for the substrates in a
diversified partition. The gap remains positive as long as each substrate's volume
share $f_i$ is small enough to compensate.

Two instructive limits bracket the regime of interest:
\begin{itemize}
\item \textbf{Full-saturation contagion} ($c_i = c = 1$): the
    condition reduces to $1 > f_{\max}$, which is exactly the
    substrate-capability-non-redundancy assumption from
    \citet{lasser2026horizon}. No additional structure is required.
\item \textbf{Unbalanced partition under contagion amplification}
    ($c_i > c$, $f_{\max}$ large): the condition can fail. For
    example, $c = 0.5$, $c_{\max} = 1.0$, $f_{\max} = 0.6$ gives
    $0.5 \not> 0.6$ and the gap reverses. In this regime
    diversification provides no dependency-risk advantage over the
    single-substrate baseline: the adversary still targets the
    dominant substrate, and the contagion saturates it more
    thoroughly than in the homogeneous case. The trust model should
    detect this regime and avoid unbalanced partitions when contagion
    risks dominate; the bound flags the edge case rather than
    assuming it away.
\end{itemize}

The analysis assumes cascade confinement: $c_i$ measures only the
fraction of substrate $\mathcal{S}_i$ lost when $\mathcal{S}_i$ is
targeted, not propagation into other substrates. Cross-substrate
cascades --- a memetic attack on $\mathcal{S}_1$ that triggers a
cascade in $\mathcal{S}_2$ through inter-substrate communication
channels --- are outside the scope of the current bound and require
a separate treatment beyond the current SCM.

\begin{corollary}[Balanced-Partition Substrate Count]
\label{cor:m_bound}
Under a balanced substrate partition with $f_i = 1/m$ for all $i$,
the adversarial balance condition of
Proposition~\ref{prop:expected_risk} is equivalent to
\begin{equation}
\label{eq:m_bound}
m > \frac{c_{\max}}{c},
\end{equation}
where $c_{\max} = \max_i c_i$ is the worst-case per-substrate
cascade fraction.
\end{corollary}

\begin{proof}
With $f_i = 1/m$, the adversarial balance condition
$c > \max_i (f_i \cdot c_i)$ reduces to
$c > \max_i (c_i / m) = c_{\max} / m$, equivalent to Equation~\eqref{eq:m_bound}.
\end{proof}

Corollary~\ref{cor:m_bound} sharpens the minimax bound of
\citet{lasser2026horizon}, which establishes only that $m \geq 2$
strictly dominates $m = 1$. The minimax statement is the special
case $c_{\max} = c$ of Equation~\eqref{eq:m_bound}: cascade
intensity unchanged by partition size gives $m > 1$, recovering
$m \geq 2$. Weakening either parameter raises the required
substrate count:
\begin{itemize}
\item \textbf{Moderate contagion} ($c = 0.5$, $c_{\max} = 1.0$):
    $m > 2$, so $m \geq 3$. Binary diversification is insufficient.
\item \textbf{High contagion} ($c = 0.1$, $c_{\max} = 1.0$):
    $m > 10$, so $m \geq 11$. The substrate-count requirement grows
    inversely with the homogeneous cascade fraction.
\item \textbf{Low contagion} ($c \approx c_{\max} \approx 0.3$,
    localized failure modes): $m > 1$, recovering the minimax bound.
\end{itemize}
Practically: the minimum number of substrates required for
diversification to strictly help depends on the measured contagion
intensity, not just on the structural claim that $m \geq 2$.
A binary partition is correct for threat classes where the
homogeneous cascade already saturates ($c \to 1$), but
underprovisions for contagion-driven threats where smaller
substrates saturate faster than the homogeneous baseline. For
unbalanced partitions, the operational form
$f_{\max} < c / c_{\max}$ is the direct test: it bounds the largest
substrate's share rather than specifying a count, and is the
condition the trust model should check when assessing a given
partition's dependency-risk posture.

\subsection{Risk-Adjusted Trajectory Comparison}
\label{sec:risk_adjusted}

The per-step expected risk gap compounds at the same rate as growth, unlike
the one-time minimax penalty in \citet{lasser2026horizon}. Under the
assumption that $\Delta_{\mathrm{risk}}$ is stationary across the
trajectory --- that is, $p_{\mathrm{adv}}$, the per-substrate shares
$f_i$, and the cascade fractions $c_i$ are time-invariant in
expectation --- the geometric series sums to the standard discounted
form:
\begin{equation}
\label{eq:risk_adjusted_per_step}
\Vdisc^{\mathrm{div},\mathrm{risk}} - \Vdisc^{D,\mathrm{risk}}
    = \frac{r_{\mathrm{ext}} + \Delta_{\mathrm{risk}}}{1 - \gamma} - \Delta_0
\end{equation}
where $\Delta_{\mathrm{risk}} =
\mathbb{E}[\Risk_{\mathrm{dep}}^{\mathrm{single}}] -
\mathbb{E}[\Risk_{\mathrm{dep}}^{\mathrm{mixed}}] > 0$ is the per-step
risk gap from Proposition~\ref{prop:expected_risk}. The threshold for
diversity dominance becomes:
\begin{equation}
\label{eq:risk_adjusted_threshold}
\gamma^*_{\mathrm{risk}} = 1 - \frac{r_{\mathrm{ext}} +
    \Delta_{\mathrm{risk}}}{\Delta_0}
\end{equation}
The risk gap $\Delta_{\mathrm{risk}}$ enters additively with
$r_{\mathrm{ext}}$, effectively increasing the external contribution by the
per-step risk mitigation value of cross-substrate diversity. When
the risk gap is non-stationary --- e.g., the threat landscape
evolves, or the substrate partition itself changes during the
trajectory --- the closed-form $(1 - \gamma)^{-1}$ denominator
becomes an approximation, and the exact value is
$\sum_{t=0}^{\infty} \gamma^t \cdot \Delta_{\mathrm{risk}}(t)$
computed from a time-varying risk profile. The trust model's
standard non-stationarity assumption (agents and threats evolve on
timescales slower than the discount horizon) makes the stationary
approximation adequate for planning; rapidly-varying threat
landscapes require a finer-grained value-iteration calculation
instead of the closed form.

\paragraph{Quantitative impact.} For a population with $p_{\mathrm{adv}} =
0.01$ (1\% per-step probability of substrate-targeting event) and cascade
fraction $c = 0.5$ (half the targeted substrate's capabilities are lost in
a typical event), $\Delta_{\mathrm{risk}} = 0.01 \times \Delta_{\mathrm{div}}
\times 0.5$ where $\Delta_{\mathrm{div}}$ is the substrate diversification
advantage from Proposition~4 of \citet{lasser2026horizon}. If
$\Delta_{\mathrm{div}}$ is 10\% of $\volL(G)$, then $\Delta_{\mathrm{risk}}
\approx 0.0005 \times \volL(G)$: the per-step risk cost of single-substrate
concentration is 0.05\% of total volume per step. Over the full planning
horizon ($1/(1-\gamma)$ steps), this compounds to $0.05\% \times
\volL(G) / (1-\gamma)$, which for $\gamma = 0.99$ is $5\% \times
\volL(G)$---substantial.

\paragraph{Sensitivity analysis.} Table~\ref{tab:sensitivity} shows
how $\Delta_{\mathrm{risk}} / \volL(G)$ varies over plausible
parameter ranges. The per-step risk gap scales linearly with
$p_{\mathrm{adv}}$ and $c$, and the compounded effect over the
planning horizon is $\Delta_{\mathrm{risk}} / (1 - \gamma)$.
The 80\% upper-bound entry corresponds to a 5\% adversarial rate
with high contagion ($c = 0.8$) and substantial divergence
($\Delta_{\mathrm{div}} = 0.20$), representing a
compromised-substrate scenario rather than a typical operating
point. A more representative regime is
$p_{\mathrm{adv}} = 0.01$, $c = 0.5$, $\Delta_{\mathrm{div}} =
0.10$, which yields a 5\% compounded gap---material but
containable by the verification architecture of
Section~\ref{sec:attention_not_action}.

\begin{table}[h]
\centering
\small
\begin{tabular}{@{}ccccc@{}}
\toprule
$p_{\mathrm{adv}}$ & $c$ & $\Delta_{\mathrm{div}}/\volL$ & $\Delta_{\mathrm{risk}}/\volL$ (per step) & Compounded at $\gamma = 0.99$ \\
\midrule
0.001 & 0.3 & 0.10 & $3 \times 10^{-5}$   & $0.3\%$ \\
0.01  & 0.3 & 0.10 & $3 \times 10^{-4}$   & $3.0\%$ \\
0.01  & 0.5 & 0.10 & $5 \times 10^{-4}$   & $5.0\%$ \\
0.01  & 0.5 & 0.20 & $1 \times 10^{-3}$   & $10.0\%$ \\
0.01  & 0.8 & 0.10 & $8 \times 10^{-4}$   & $8.0\%$ \\
0.05  & 0.5 & 0.10 & $2.5 \times 10^{-3}$ & $25.0\%$ \\
0.05  & 0.8 & 0.20 & $8 \times 10^{-3}$   & $80.0\%$ \\
\bottomrule
\end{tabular}
\caption{Sensitivity of the per-step dependency-risk gap
    $\Delta_{\mathrm{risk}} = p_{\mathrm{adv}} \cdot
    \Delta_{\mathrm{div}} \cdot c$ and its compounded effect over
    the planning horizon at $\gamma = 0.99$
    ($1/(1 - \gamma) = 100$ effective steps). The compounded
    percentage is $\Delta_{\mathrm{risk}} / [(1 - \gamma) \cdot
    \volL(G)] \times 100$.}
\label{tab:sensitivity}
\end{table}

\paragraph{Interaction with the horizon paper's cooperative-novelty result.}
The per-step risk gap $\Delta_{\mathrm{risk}}$ is amplified by the
cross-substrate cooperative capabilities analyzed in
\citet{lasser2026horizon}: any
cooperative capability requiring agents on both substrates is
destroyed when either substrate is targeted, so
$|\Delta \volL(\mathcal{S}_i)|$ includes those cooperative losses on
top of the individual-capability loss on $\mathcal{S}_i$. Under the
horizon paper's cooperative-novelty result, these cross-substrate
cooperative capabilities are the primary growth channel, so their
destruction represents a disproportionate fraction of the total
contraction during a substrate-targeting event. The cooperative
novelty and dependency-risk analyses reinforce each other: domination
is costly both because it eliminates the cross-substrate cooperative
growth channel \emph{and} because it concentrates adversarial
exposure to the single surviving substrate
\citep{lasser2026horizon}.

\paragraph{Summary.} The adversarial balance condition
$c > \max_i (f_i \cdot c_i)$ converts the companion paper's minimax
bound into an expected per-step cost $\Delta_{\mathrm{risk}}$ that
compounds as $\Delta_{\mathrm{risk}} / (1 - \gamma)$ in the
discounted value function, making dependency risk a first-class term
in the horizon calculation rather than a one-time structural
penalty. The condition is explicit about contagion dynamics and
inverts (Corollary~\ref{cor:m_bound}) into a quantitative bound on
the required substrate count: $m > c_{\max}/c$ for balanced
partitions, or $f_{\max} < c / c_{\max}$ for arbitrary partitions.
This sharpens the minimax $m \geq 2$ into a measurement-driven
requirement that scales with contagion intensity. The regime where
the condition fails --- unbalanced partitions under contagion
amplification --- is a regime where diversification genuinely does
not help, which the trust model can detect and route around.
